%! Observational Studies — Unit 1, Lesson 1.5 — Group Activity
\documentclass[10pt]{article}
\usepackage{saar-article}
\usepackage{saar-boxes}
\newcommand{\TallMath}[1]{$\displaystyle #1\rule[-1.4em]{0pt}{3.2em}$}

\begin{document}

\pageheader{Unit 1, Lesson 1.5}{Group Activity}
% NAMESTRIP: no \namedateperiod here. The student writes their name once, on the
% cover this component is stapled behind. See references/conventions.md.

\vspace{0.15in}

\begin{headlinebox}{frost}
\small
\textbf{Cancel the Study Center? --- Riverbend High School.} Riverbend has $900$
students and an after-school Study Center that any student may walk into. A
counselor pulls a random sample of $150$ student records from the full roster.
Of those $150$, \textbf{$60$ used the Study Center at least once a week} last
quarter and \textbf{$90$ never did}. Among the $60$ users, $39$ passed every
class; among the $90$ non-users, $72$ passed every class. The principal reads
the numbers and drafts a memo: \emph{``The Study Center is hurting grades. Shut
it down.''}
\end{headlinebox}

\vspace{0.03in}

\boxguard[30]
\begin{tcolorbox}[colback=white, colframe=black!40, breakable,
    title=\textbf{Tier R --- Remediate},
    fonttitle=\bfseries, arc=2mm, left=3mm, right=3mm, top=1mm, bottom=1mm]
\small
\begin{enumerate}[leftmargin=1.6em, itemsep=2pt, topsep=3pt]

  \item Name the numbers in this study.

        Population: \blank{2.6cm} \hfill Sample: \blank{1.2cm} \hfill
        Used the Center: \blank{1.2cm} \hfill Did not: \blank{1.2cm}

  \item Nobody at Riverbend told a single student to go to the Study Center.

        \textbf{(a)} Is this an observational study? \blank{2cm}

        \textbf{(b)} Who decided which group each student ended up in?
        \blank{4.6cm}

  \item Of the $60$ students who used the Study Center, $39$ passed every class.
        What percent is that?

        \begin{work}
          \dfrac{39}{60} &= 0.65 \\
                         &= 65\%
        \end{work}

  \item Of the $90$ students who never used it, $72$ passed every class. What
        percent is that?

        \begin{work}
          \dfrac{72}{90} &= 0.80 \\
                         &= 80\%
        \end{work}

  \item Which group has the higher pass rate, and by how many percentage points?
        \writeline
\end{enumerate}
\end{tcolorbox}

\vspace{0.03in}

\boxguard
\begin{tcolorbox}[colback=white, colframe=black!40, breakable,
    title=\textbf{Tier A --- Approaching Proficiency},
    fonttitle=\bfseries, arc=2mm, left=3mm, right=3mm, top=1.5mm, bottom=1.5mm]
\small
\begin{enumerate}[leftmargin=1.6em, itemsep=3pt, topsep=3pt]

  \item Put the sample together.

        \textbf{(a)} How many of the $150$ students in the sample passed every
        class, and what percent is that?

        \begin{work}
          39 + 72 &= 111 \text{ students} \\
          \dfrac{111}{150} &= 0.74 = 74\%
        \end{work}

        \textbf{(b)} About how many of all $900$ Riverbend students does that
        predict?

        \begin{work}
          0.74 \times 900 &= 666 \text{ students}
        \end{work}

  \item Name the two variables the counselor compared.

        Explanatory: \blank{4.4cm} \hfill Response: \blank{4.4cm}

  \item Write the counselor's finding as an \textbf{association} statement ---
        both percents, the gap, and no word like \emph{causes} or \emph{hurts}.
        \writelines{2}

  \item The principal's memo says the Study Center \emph{caused} the lower pass
        rate. Explain why this study cannot support that.
        \writelines{2}
\end{enumerate}
\end{tcolorbox}

\vspace{0.03in}

\boxguard
\begin{tcolorbox}[colback=white, colframe=black!40, breakable,
    title=\textbf{Tier E --- Extension},
    fonttitle=\bfseries, arc=2mm, left=3mm, right=3mm, top=2mm, bottom=2mm]
\small
\begin{enumerate}[leftmargin=1.6em, itemsep=5pt, topsep=3pt]

  \item The counselor checks one more column: how each student did the
        \emph{previous} quarter, before the Study Center opened. Of the $60$
        users, $45$ had failed at least one class. Of the $90$ non-users, only
        $9$ had.

        \textbf{(a)} Users who had already failed a class: \blank{2cm} \quad
        Non-users: \blank{2cm}

        \textbf{(b)} What does that tell you about \emph{who} walks into the
        Study Center?
        \writeline

  \item Does the fact you found in item 1 make the Study Center look
        \emph{better} or \emph{worse} than it really is? Explain how you know.
        \writelines{2}

  \item The principal wants a study that would actually answer ``does the Study
        Center help?'' For each option, say what it fixes and what it costs.

        \textbf{(a)} Compare each user's grades this quarter to that same
        student's grades last quarter.
        \writeline

        \textbf{(b)} Assign $150$ randomly chosen students to attend, and $150$
        randomly chosen students not to.
        \writeline
\end{enumerate}
\end{tcolorbox}

\end{document}
